\documentclass[12pt,largemargins]{quiz}

\usepackage{tikz}
\usepackage{graphicx}
\usetikzlibrary{automata, positioning, arrows.meta, calc, fit}

% --- Quiz Metadata ---
\newcommand{\eventclass}{نظریه زبان‌ها و ماشین‌ها}
\newcommand{\eventtype}{آزمون کوتاه}           % یا: کوییز شماره
\newcommand{\eventnum}{۱}
\newcommand{\eventname}{آتوماتای متناهی قطعی (DFA)}
\newcommand{\eventcntrbs}{استاد: دکتر فدایی}
\newcommand{\eventdatetype}{تاریخ برگزاری:}
\newcommand{\eventdate}{۲۵ آبان ۱۴۰۵}
\newcommand{\quiztime}{۱۵ دقیقه}

\begin{document}
\maketitle

\question
یک DFA طراحی کنید که الفبای آن $\Sigma = \{0, 1\}$ است و تمام رشته‌هایی را می‌پذیرد که تعداد رقم‌های ۱ در آن‌ها فرد باشد.

\answerbox{3.5cm}

% در صورت نیاز به نوشتن پاسخ تشریحی:
% \solution
% پاسخ تشریحی در این بخش قرار می‌گیرد.

\question*{سؤال دوم: ویژگی‌های زبان‌های منظم}
با فرض زبان $L = \{w \in \{a, b\}^* \mid |w| \ge 2\}$، به بخش‌های زیر پاسخ دهید:

\begin{alphaparts}
  \questionpart آیا این زبان منظم است؟ یک عبارت منظم معادل برای آن بنویسید.
  \answerbox{2.5cm}
  
  \questionpart گرامر مستقل از متن یا خطی متناظر با این زبان را بنویسید.
  \answerbox{2.5cm}
\end{alphaparts}

\end{document}

% To compile via command line:
% latexmk QUIZ.tex
% Or:
% xelatex -synctex=1 -interaction=nonstopmode -shell-escape QUIZ.tex
